% Lösungen Einheit 2 von 5 – Prozentsatz berechnen (zusammenbau v0.3)
\einheitenkopf{Einheit 2 von 5 · Prozentsatz berechnen}

\erg{12}{a) $28$ Kinder}
\erg{13}{a) $10$ gleiche Teile}
\erg{14}{a) $20\,\%$ \quad b) $17\,\%$ \quad c) $\frac{19}{50} = \frac{38}{100} = 38\,\%$ \quad d) $27 : 60 = 0{,}45 = 45\,\%$ \quad e) $17 : 24 = 0{,}7083\ldots \approx 70{,}8\,\%$ \quad f) Ganzes $20$; $6 : 20 = 0{,}3 = 30\,\%$ \quad g) Rabatt $20\,€$; $20 : 80 = 0{,}25 = 25\,\%$ \quad h) Ganzes $24$ Lose; $3 : 24 = 0{,}125 = 12{,}5\,\%$ \quad i) $1\,080 : 40\,000 = 0{,}027 = 2{,}7\,\%$ \quad j) $6$ von $13$ sind größer; $6 : 13 \approx 0{,}462 \approx 46{,}2\,\%$}
\erg{15}{a) z. B. $3$ von $10$ oder $15$ von $50$ (Teil : Ganzes $= 0{,}3$)}
\erg{16}{a) Fehler: Ganzes durch Teil gerechnet. Richtig: $12 : 48 = 0{,}25 = 25\,\%$. \quad b) Gefragt ist, welcher Anteil vom Ganzen der Teil ist – das Ganze sind die $100\,\%$, also steht es im Nenner. \quad c) $25\,\%$ ($1$ von $4$ Teilen) \quad d) übrig $3$ GB, $3 : 20 = 0{,}15 = 15\,\%$; pro Tag $17 : 20 = 0{,}85$ GB, für $10$ Tage bräuchte man $10 \cdot 0{,}85 = 8{,}5$ GB – es reicht nicht.}
